
\section*{3. Identity component of a group locally of finite
  presentation}
\phantomsection\label{sec:sga3-VIB-3}
\addcontentsline{toc}{section}{3. Identity component of a group
  locally of finite presentation}

\medskip
\phantomsection\label{par:sga3-VIB-3.0}%
\noindent\textbf{3.0.---}
Given a subset \(A\) (respectively, \(B\)) of an \(S\)-scheme
\(X\) (respectively, \(Y\)), by abuse of notation
\[
  A\times_S B
\]
will denote the subset of \(X\times_S Y\) consisting of the points
whose first projection belongs to \(A\) and whose second projection
belongs to \(B\).

Given a subset \(A\) of an \(S\)-group \(G\), we shall say that
\(A\) is \emph{stable under the group law of \(G\)} if
\[
  c(A)\subset A
  \qquad\text{and}\qquad
  \mu(A\times_S A)\subset A.
\]

\medskip
\phantomsection\label{def:sga3-VIB-3.1}%
\noindent\textbf{Definition 3.1.---}
Let \(G\) be an \(S\)-functor in groups satisfying the following
condition:
\[
  \tag{+}\label{sga3:vi:definition:3-1:equation:plus}
  \text{for every \(s\in S\), the functor
  \(G_s=G\otimes_S\kappa(s)\) is representable.}
\]
We then denote by \(\underline{G_s^0}\) the connected component of
the identity element of the \(\kappa(s)\)-group \(G_s\).%
\footnote{\label{note:sga3-VIB-22}Editor's note: this is a slightly
abusive notation, but it is compatible with the notation of
Expos\'e~VI\(_A\), \SGARefLink{cor:sga3-VIB-2.3}{2.3} when this connected component is the
underlying topological space of an open group subscheme \(G^0\) of
\(G\); cf.\ Expos\'e~VI\(_A\), \SGARefLink{par:sga3-VIA-2.2bis}{2.2.b}is.}
We define a sub-\(S\)-functor in groups of \(G\), called the
\emph{identity component of \(G\)} and denoted by \(G^0\), by
setting, for every morphism \(T\to S\),
\[
  G^0(T)
    =
  \bigl\{u\in G(T)\bigm|
    \text{for every \(s\in S\), }
    u_s(T_s)\subset\underline{G_s^0}\bigr\}.
\]
Thus we have defined the functor
\[
  G\longmapsto G^0
\]
from
\[
  \bigl(\widehat{\mathbf{Sch}}_{/S}\bigr)\text{-gr.}
  \quad\text{to}\quad
  \bigl(\widehat{\mathbf{Sch}}_{/S}\bigr)\text{-gr.}
\]

\medskip
\phantomsection\label{rem:sga3-VIB-3.2}%
\noindent\textbf{Remark 3.2.---}
\textup{(i)} Let \(G\) be an \(S\)-functor in groups satisfying
condition \((+)\). Then
\[
  \operatorname{Lie}(G^0/S)=\operatorname{Lie}(G/S)
\]
by Expos\'e~II.%
\footnote{\label{note:sga3-VIB-23}Editor's note: we have added what
follows.}
Indeed, for every \(S\)-scheme \(T\), denote by \(I_T\) the
\(T\)-scheme of dual numbers over \(T\), and by
\[
  \tau:T\longrightarrow I_T
\]
the zero section (cf.\ \SGARefLink{sga3:ii:definition:2-1}{Expos\'e~II, 2.1}). By definition,
\[
  \operatorname{Lie}(G/S)(T)
    =
  \{u\in G(I_T)\mid u\circ\tau=e\},
\]
where \(e:T\to G\) denotes the composite of \(T\to S\) with the
identity section \(S\to G\). Likewise,
\[
  \begin{aligned}
  \operatorname{Lie}(G^0/S)(T)
    &=
  \{u\in G^0(I_T)\mid u\circ\tau=e\} \\
    &=
  \bigl\{u\in G(I_T)\bigm|
    u\circ\tau=e
    \ \text{and}\
    u_s\bigl((I_T)_s\bigr)\subset\underline{G_s^0}
    \text{ for every \(s\in S\)}\bigr\}.
  \end{aligned}
\]
For every \(s\in S\), the schemes \(T_s\) and
\[
  (I_T)_s=I_{T_s}
\]
have the same underlying set. Hence, if
\(u\in\operatorname{Lie}(G/S)(T)\), then
\[
  u_s(I_{T_s})=e_s,
\]
where \(e_s\) denotes the identity point of \(G_s\), and therefore
\(u\in\operatorname{Lie}(G^0/S)(T)\). Thus the inclusion
\[
  \operatorname{Lie}(G^0/S)(T)
    \subset
  \operatorname{Lie}(G/S)(T)
\]
is an equality for every \(T\), and consequently
\[
  \operatorname{Lie}(G^0/S)=\operatorname{Lie}(G/S).
\]

\textup{(ii)} Let \(G\) and \(G'\) be two \(S\)-functors in groups
satisfying \((+)\). Then:
\begin{enumerate}
\item[\textup{a)}] if \(G\subset G'\), then \(G^0\subset G'^0\);
\item[\textup{b)}] if \(G\subset G'\) and \(G'^0\subset G\), then
  \(G^0=G'^0\);
\item[\textup{c)}] if, for every \(s\in S\), \(G_s\) is locally of
  finite type over \(\kappa(s)\), then \(G^0\) satisfies condition
  \((+)\), by Expos\'e~VI\(_A\), \SGARefLink{cor:sga3-VIB-2.3}{2.3}, and
  \[
    (G^0)^0=G^0.
  \]
\end{enumerate}

\medskip
\phantomsection\label{prop:sga3-VIB-3.3}%
\noindent\textbf{Proposition 3.3.---}
\emph{Let \(G\) be an \(S\)-functor in groups satisfying condition
\((+)\), and let \(S'\) be an \(S\)-scheme. Then}
\[
  (G\times_S S')^0=G^0\times_S S'.
\]
\emph{In other words, the functor \(G\mapsto G^0\) commutes with
base change; that is, the following diagram is commutative:}
\begin{center}
  \phantomsection\label{diag:sga3-VIB-3.3-base-change}%
  \begin{tikzcd}[row sep=large,column sep=huge]
    \bigl(\widehat{\mathrm{Sch}}_{/S}\bigr)\text{-}\mathrm{gr.}
      \arrow[r,"{G\mapsto G^0}"]
      \arrow[d] &
    \bigl(\widehat{\mathrm{Sch}}_{/S}\bigr)\text{-}\mathrm{gr.}
      \arrow[d] \\
    \bigl(\widehat{\mathrm{Sch}}_{/S'}\bigr)\text{-}\mathrm{gr.}
      \arrow[r] &
    \bigl(\widehat{\mathrm{Sch}}_{/S'}\bigr)\text{-}\mathrm{gr.}
  \end{tikzcd}
\end{center}

It is enough to check that, for every \(s'\in S'\), whose image in
\(S\) we denote by \(s\),
\[
  \begin{aligned}
  \bigl((G\times_S S')\otimes_{S'}\kappa(s')\bigr)^0
    &=
  \bigl(G_s\otimes_{\kappa(s)}\kappa(s')\bigr)^0 \\
    &=
  G_s^0\otimes_{\kappa(s)}\kappa(s').
  \end{aligned}
\]
This follows from Expos\'e~VI\(_A\), \SGARefLink{rem:sga3-VIB-2.1.2}{2.1.2}. Note, for later use in
\SGARefLink{prop:sga3-VIB-4.2}{4.2}, that we have not used the group structure of \(G_s\), but
only the fact that \(G_s^0\) has a rational point, namely
\(\epsilon(s)\), and is therefore geometrically connected
(see also EGA~IV\(_2\), 4.5.14).

\medskip
\phantomsection\label{case:sga3-VIB-3.4}%
\noindent\textbf{Special Case 3.4.---}
Let \(G\) be an \(S\)-group scheme, and denote by
\(\underline{G^0}\) the subset of \(G\) equal to the union of the
\(\underline{G_s^0}\) as \(s\) runs through \(S\). Then
\(\underline{G^0}\) is a subset of \(G\) stable under the group law
of \(G\) (cf.\ \hyperref[par:sga3-VIB-3.0]{3.0}), and, for every morphism \(S'\to S\),
\[
  G^0(S')
    =
  \bigl\{u\in G(S')\bigm|
    u(S')\subset\underline{G^0}\bigr\}.
\]
When \(\underline{G^0}\) is an open subset of \(G\), it is clear
that \(G^0\) is represented by the subscheme of \(G\) induced on
the open subset \(\underline{G^0}\).

\medskip
\phantomsection\label{prop:sga3-VIB-3.5}%
\noindent\textbf{Proposition 3.5.---}
\emph{Let \(S\) be a quasi-compact and quasi-separated scheme, and
let \(G\) be an \(S\)-group with fibers locally of finite type.
Then there is a quasi-compact open subset \(U\) of \(G\) containing
\(\underline{G^0}\).}

Since the identity section \(\epsilon\) is an immersion,
\(\epsilon(S)\) is a quasi-compact subspace of \(G\). Thus there is
a quasi-compact open subset \(V\) of \(G\) containing
\(\epsilon(S)\). Since \(S\) is quasi-separated and \(V\) is
quasi-compact, \(V\) is quasi-compact over \(S\)
(EGA~IV\(_1\), 1.2.4). Hence \(V\times_SV\) is quasi-compact over
\(S\), and therefore quasi-compact. It follows that
\[
  V\cdot V=\mu(V\times_SV)
\]
is quasi-compact. Set
\[
  V_s=V\cap G_s,
  \qquad
  V_s^0=V\cap\underline{G_s^0}.
\]
Then \(V_s^0\) is open in \(\underline{G_s^0}\), and is dense in
\(\underline{G_s^0}\) because \(G_s^0\) is irreducible
(Expos\'e~VI\(_A\), \SGARefLink{cor:sga3-VIB-2.4}{2.4}). Therefore
\[
  V_s^0\cdot V_s^0=\underline{G_s^0}
  \qquad
  \text{(Expos\'e~VI\(_A\), 0.5),}
\]
which shows that
\[
  V_s\cdot V_s\supset\underline{G_s^0},
  \qquad\text{and hence}\qquad
  V\cdot V\supset\underline{G^0}.
\]
Finally, since \(V\cdot V\) is quasi-compact, there is a
quasi-compact open subset \(U\) of \(G\) containing \(V\cdot V\),
and therefore containing \(\underline{G^0}\).

\medskip
\phantomsection\label{cor:sga3-VIB-3.6}%
\noindent\textbf{Corollary 3.6.---}
\emph{Let \(G\) be an \(S\)-group whose fibers are locally of
finite type and connected. Then \(G\) is quasi-compact over \(S\).}

The assertion is local on \(S\) (EGA~I, 6.6.1), so we are reduced
to the case where \(S\) is affine. By 3.5, there is then a
quasi-compact open subset \(U\) of \(G\) containing
\[
  \underline{G^0}=G.
\]
Thus \(G\) is quasi-compact, and hence quasi-compact over the affine
scheme \(S\) (EGA~I, 6.6.4(v)).
